% The credential data model (§4.1): rules and guarantees that are in Lean.
% Written by hand.  Do not write \leanok; scripts/render_status.py adds it.
% The candidates of this area that have no Lean yet are in
% credentials-candidates.tex.

\section{Credential data model}

\begin{definition}[Data elements and their names (C-25)]
  \label{def:element}
  \lean{Rules.OneId.Mdoc.element?, Rules.OneId.Mdoc.wellFormed, Rules.OneId.Mdoc.nameSpace?, Rules.OneId.NameSpace.element?}
  \uses{def:mdoc}
  Each data element of a credential is identified by the name of its
  namespace and its identifier.  The names of a credential are unique when no
  two namespaces share a name and no two data elements of one namespace share
  an identifier.  The guide leaves the structure to ISO/IEC 18013-5, which
  this project has not read, and does not say that names are unique.

  External: ISO/IEC 18013-5:2021 §6.3.1, §8.3.1, §8.3.2.

  Source: \href{https://standards.atlassian.net/wiki/spaces/OneID/pages/2947416082/Technical+Implementation+Guide}{One ID guide} §4.1, §4.1.1.
\end{definition}

\begin{theorem}[Every data element can be requested on its own]
  \label{thm:element-addressable}
  \lean{Rules.OneId.element_addressable}
  \uses{def:element}
  In a credential whose names are unique, every data element is found by the
  name of its namespace and its identifier.
\end{theorem}
\begin{proof}
  \uses{def:element}
  With unique names, the search for a name finds the element that carries it.
\end{proof}

\begin{theorem}[No value for a name the credential lacks]
  \label{thm:element-absent}
  \lean{Rules.OneId.element_absent}
  \uses{def:element}
  A credential returns no value for a namespace name and an identifier when
  none of its namespaces of that name holds a data element with that
  identifier.
\end{theorem}
\begin{proof}
  \uses{def:element}
  The search for the namespace or for the element finds nothing.
\end{proof}

\begin{theorem}[The names in a credential are unique]
  \label{thm:names-unique}
  \lean{Rules.OneId.photoid_names_unique, Rules.OneId.boardingPass_names_unique}
  \uses{def:element, def:photoid, def:boarding-pass}
  The names in an accepted photoid, and in an accepted boarding pass, are
  unique.  Not proven.

  Counterexample: a photoid whose base identity namespace holds two family
  names, SMITH and JONES, is accepted under the reading that names may
  repeat.  A relying party that requests the family name receives one of the
  two.  A boarding pass with two boarding pass strings is accepted in the same
  way.

  Open point: the guide does not say that namespace names and data element
  identifiers are unique.  May a name repeat within a credential, and is such
  a credential rejected?

  Lookup: ISO/IEC 18013-5:2021 §8.3.1, §8.3.2.

  Source: \href{https://standards.atlassian.net/wiki/spaces/OneID/pages/2947416082/Technical+Implementation+Guide}{One ID guide} §4.1.
\end{theorem}

\begin{definition}[Photoid requirements (C-27)]
  \label{def:photoid}
  \lean{Rules.OneId.Photoid.violations, Rules.OneId.PhotoidReading, Rules.OneId.PhotoidTables, Rules.OneId.ElementTable, Rules.OneId.TableEntry, Rules.OneId.Presence, Rules.OneId.Encoding}
  \uses{def:element, def:finding}
  A credential is accepted as a photoid only when it has the document type
  \texttt{org.iso.23220.photoid.1} and the base identity namespace
  \texttt{org.iso.23220.1}, with every data element Table C.1 of ISO/IEC TS
  23220-4 marks mandatory and each listed element in the encoding the table
  gives.  The tables are a parameter, since this project has not read them.
  The points the guide leaves undecided are a reading: whether the portrait
  is mandatory, whether data the guide does not list is allowed, whether the
  portrait must be the same as data group 2, whether a recommended namespace
  must be complete once it is present, and whether names may repeat.

  External: ISO/IEC TS 23220-4:2026 Annex C, §C.2.1 and Tables C.1 to C.3;
  ISO/IEC TS 23220-2:2026.

  Source: \href{https://standards.atlassian.net/wiki/spaces/OneID/pages/2947416082/Technical+Implementation+Guide}{One ID guide} §4.1.1, §2.
\end{definition}

\begin{theorem}[A photoid has the photoid document type]
  \label{thm:photoid-doctype}
  \lean{Rules.OneId.photoid_docType}
  \uses{def:photoid}
  Under every reading and every content of the tables, an accepted photoid
  has the photoid document type.
\end{theorem}
\begin{proof}
  \uses{def:photoid}
  Another document type is a finding.
\end{proof}

\begin{theorem}[A photoid has the base identity namespace]
  \label{thm:photoid-has-base}
  \lean{Rules.OneId.photoid_has_base}
  \uses{def:photoid}
  Under every reading and every content of the tables, an accepted photoid
  has the base identity namespace.
\end{theorem}
\begin{proof}
  \uses{def:photoid}
  A missing namespace is a finding.
\end{proof}

\begin{theorem}[A photoid has every mandatory base element]
  \label{thm:photoid-mandatory-base}
  \lean{Rules.OneId.photoid_has_mandatory_base}
  \uses{def:photoid}
  An accepted photoid holds every data element Table C.1 marks mandatory in
  the base identity namespace, whatever the table lists.
\end{theorem}
\begin{proof}
  \uses{def:photoid}
  Each mandatory element that is missing is a finding.
\end{proof}

\begin{theorem}[Base elements have the encoding of their table]
  \label{thm:photoid-base-encoding}
  \lean{Rules.OneId.photoid_base_encoding}
  \uses{def:photoid}
  In an accepted photoid, every data element of the base identity namespace
  that Table C.1 lists has the encoding the table gives it, whatever the
  table says.  The model tells a text string, a byte string and a yes or no
  value apart.
\end{theorem}
\begin{proof}
  \uses{def:photoid}
  A listed element in another encoding is a finding.
\end{proof}

\begin{theorem}[Every photoid has a portrait]
  \label{thm:photoid-has-portrait}
  \lean{Rules.OneId.photoid_has_portrait}
  \uses{def:photoid, def:c04, def:c103}
  Every accepted photoid has a portrait.  Not proven.

  Counterexample: a photoid with a family name and a given name and no
  portrait is accepted when Table C.1 does not mark the portrait mandatory and the
  guide's sentence is read as a rule about the portrait's content only.

  Open point: the guide says the portrait MUST hold the passport's image and
  that the business requirements need it, and it also refers the presence of
  every base element to Table C.1.  Is the portrait mandatory in a One ID
  photoid, whatever Table C.1 says?

  Lookup: ISO/IEC TS 23220-4:2026 Annex C, Table C.1.  If the table marks the
  portrait mandatory, the guarantee follows.

  Source: \href{https://standards.atlassian.net/wiki/spaces/OneID/pages/2947416082/Technical+Implementation+Guide}{One ID guide} §4.1.1.
\end{theorem}

\begin{theorem}[A photoid carries only what the guide lists]
  \label{thm:photoid-only-listed}
  \lean{Rules.OneId.photoid_only_listed}
  \uses{def:photoid}
  An accepted photoid carries only the three namespaces the guide lists, and
  in each only data elements its table lists.  Not proven.

  Counterexample: a photoid that carries a fourth namespace with a frequent
  flyer number, and a data element in its base identity namespace that Table
  C.1 does not list, is accepted under the reading that the lists are not
  closed.

  Open point: the guide lists three namespaces and refers to a table for
  each, and does not say the lists are closed.  May a One ID photoid carry a
  namespace or a data element the guide and its tables do not list?

  Source: \href{https://standards.atlassian.net/wiki/spaces/OneID/pages/2947416082/Technical+Implementation+Guide}{One ID guide} §4.1.1.
\end{theorem}

\begin{theorem}[The portrait is the same as data group 2]
  \label{thm:portrait-is-dg2}
  \lean{Rules.OneId.portrait_is_dg2}
  \uses{def:photoid}
  When an accepted photoid carries both a portrait and a data group 2, they
  are the same.  Not proven.

  Counterexample: a photoid whose portrait is one face image and whose data
  group 2 is another is accepted under the reading that ``equivalent'' is
  left to the issuer.

  Open point: the guide says the portrait MUST hold ``the dg2 or equivalent
  image data'' and does not say what equivalent means or who checks it.  When
  a photoid carries data group 2, must the portrait be that image, and must a
  relying party check it?

  Source: \href{https://standards.atlassian.net/wiki/spaces/OneID/pages/2947416082/Technical+Implementation+Guide}{One ID guide} §4.1.1.
\end{theorem}

\begin{theorem}[A photoid namespace that is present is complete]
  \label{thm:photoid-namespace-complete}
  \lean{Rules.OneId.photoid_nameSpace_complete}
  \uses{def:photoid, def:c103}
  When an accepted photoid carries the photoid namespace, the namespace holds
  every data element Table C.2 marks mandatory.  Not proven.

  Counterexample: a photoid whose photoid namespace holds only the image
  captured at enrolment is accepted under the reading that the table is only
  recommended, though the table marks a document number mandatory.

  Open point: the guide recommends ``this namespace and data elements
  presence'' in one sentence.  Once the photoid namespace is present, are the
  data elements Table C.2 marks mandatory required?

  Source: \href{https://standards.atlassian.net/wiki/spaces/OneID/pages/2947416082/Technical+Implementation+Guide}{One ID guide} §4.1.1.
\end{theorem}

\begin{theorem}[A data groups namespace that is present is complete]
  \label{thm:datagroups-namespace-complete}
  \lean{Rules.OneId.dataGroups_nameSpace_complete}
  \uses{def:photoid, def:c103}
  When an accepted photoid carries the data groups namespace, the namespace
  holds every data element Table C.3 marks mandatory.  Not proven.

  Counterexample: a photoid whose data groups namespace holds only data group
  2 is accepted under the reading that the table is only recommended, though
  the table marks the passport's security object mandatory.

  Open point: the guide recommends the data groups namespace and the presence
  of its data elements as in Table C.3 in one sentence.  Once the data groups
  namespace is present, are the data elements Table C.3 marks mandatory
  required?

  Source: \href{https://standards.atlassian.net/wiki/spaces/OneID/pages/2947416082/Technical+Implementation+Guide}{One ID guide} §4.1.1.
\end{theorem}

\begin{definition}[Source document of a photoid (C-06)]
  \label{def:photoid-source}
  \lean{Rules.OneId.Photoid.sourceFindings, Rules.OneId.SourceDocument}
  \uses{def:photoid}
  The kind of identity document a photoid was derived from is a fact the
  issuer knows.  Whether a source other than a passport is a reason to reject
  the photoid is a reading: the terminology allows an equivalent identity
  document, and §4.1.1 says the portrait MUST hold image data sourced from
  the passport.

  Source: \href{https://standards.atlassian.net/wiki/spaces/OneID/pages/2947416082/Technical+Implementation+Guide}{One ID guide} §2, §4.1.1.
\end{definition}

\begin{theorem}[A photoid comes from a passport]
  \label{thm:photoid-from-passport}
  \lean{Rules.OneId.photoid_from_passport}
  \uses{def:photoid-source}
  A photoid the system accepts was derived from a passport.  Not proven.

  Counterexample: a photoid derived from a national identity card is accepted
  under the reading that an equivalent identity document is allowed.  With
  the photoid namespace absent, nothing in the credential says which document
  it came from.

  Open point: the terminology says a photoid can be derived from a passport
  or an equivalent identity document, and §4.1.1 says its values should
  always come from the passport and its portrait MUST hold image data sourced
  from the passport.  May a One ID photoid be derived from a document other
  than a passport?

  Source: \href{https://standards.atlassian.net/wiki/spaces/OneID/pages/2947416082/Technical+Implementation+Guide}{One ID guide} §2, §4.1.1.
\end{theorem}

\begin{definition}[Photoid recommendations (C-32)]
  \label{def:photoid-advisories}
  \lean{Rules.OneId.Photoid.advisories}
  \uses{def:photoid}
  A photoid that meets every requirement is accepted, and is told which
  recommendations of the guide it does not follow: the photoid namespace with
  the travel document's data, the image captured at enrolment, the data
  groups namespace with the passport chip's data, and values derived from a
  passport.

  Open point: without the photoid namespace and the data groups namespace,
  nothing in a photoid says which document it was derived from, and a relying
  party cannot check the portrait against the passport chip.  Is that
  acceptable for the admissibility use case?

  External: ISO/IEC TS 23220-4:2026 Annex C, Tables C.2 and C.3; ICAO Doc 9303
  Part 10.

  Source: \href{https://standards.atlassian.net/wiki/spaces/OneID/pages/2947416082/Technical+Implementation+Guide}{One ID guide} §4.1.1, §2.
\end{definition}

\begin{theorem}[A missing photoid namespace is reported]
  \label{thm:advises-photoid-namespace}
  \lean{Rules.OneId.advises_photoid_nameSpace}
  \uses{def:photoid-advisories}
  A photoid that lacks the photoid namespace is told so, as a recommendation
  it does not follow.
\end{theorem}
\begin{proof}
  \uses{def:photoid-advisories}
  By the definition.
\end{proof}

\begin{theorem}[A missing data groups namespace is reported]
  \label{thm:advises-datagroups-namespace}
  \lean{Rules.OneId.advises_dataGroups_nameSpace}
  \uses{def:photoid-advisories}
  A photoid that lacks the data groups namespace is told so, as a
  recommendation it does not follow.
\end{theorem}
\begin{proof}
  \uses{def:photoid-advisories}
  By the definition.
\end{proof}

\begin{theorem}[A missing enrolment image is reported]
  \label{thm:advises-enrolment-portrait}
  \lean{Rules.OneId.advises_enrolment_portrait}
  \uses{def:photoid-advisories}
  A photoid whose photoid namespace lacks the image captured at enrolment is
  told so, whatever Table C.2 says about that image.
\end{theorem}
\begin{proof}
  \uses{def:photoid-advisories}
  By cases on whether the table marks the image mandatory.
\end{proof}

\begin{theorem}[A source other than a passport is reported]
  \label{thm:advises-source-document}
  \lean{Rules.OneId.advises_source_document}
  \uses{def:photoid-advisories}
  A photoid derived from a document other than a passport is told so.
\end{theorem}
\begin{proof}
  \uses{def:photoid-advisories}
  By the definition.
\end{proof}

\begin{theorem}[A missing recommended namespace is not a reason to reject]
  \label{thm:absent-recommended-namespace}
  \lean{Rules.OneId.recommended_nameSpace_never_required}
  \uses{def:photoid-advisories, def:photoid}
  No photoid is rejected for lacking the photoid namespace or the data groups
  namespace: neither is among the requirements a photoid can break, under
  every reading and every content of the tables.
\end{theorem}
\begin{proof}
  \uses{def:photoid}
  No part of the check reports either namespace as missing.
\end{proof}

\begin{definition}[Boarding pass requirements (C-36)]
  \label{def:boarding-pass}
  \lean{Rules.OneId.BoardingPass.violations, Rules.OneId.BoardingPassReading, Rules.OneId.Resolution792, Rules.OneId.BoardingPass.bcbpString}
  \uses{def:element, def:finding}
  A credential is accepted as a boarding pass only when it has the document
  type \texttt{org.iata.boardingpass.1}, unique names, and the namespace
  \texttt{org.iata.bcbp.1} with the data element \texttt{bcbp\_string}, a text
  string that follows IATA Resolution 792.  What the resolution decides is a
  parameter, since this project has not read it.  Whether a boarding pass may
  carry anything else, and whether names may repeat, is a reading.

  External: IATA Resolution 792; IETF RFC 8949.

  Source: \href{https://standards.atlassian.net/wiki/spaces/OneID/pages/2947416082/Technical+Implementation+Guide}{One ID guide} §4.1.2.
\end{definition}

\begin{theorem}[A boarding pass has the boarding pass document type]
  \label{thm:bp-doctype}
  \lean{Rules.OneId.boardingPass_docType}
  \uses{def:boarding-pass}
  Under every reading, an accepted boarding pass has the boarding pass
  document type.
\end{theorem}
\begin{proof}
  \uses{def:boarding-pass}
  Another document type is a finding.
\end{proof}

\begin{theorem}[A boarding pass holds one boarding pass string]
  \label{thm:bp-has-string}
  \lean{Rules.OneId.boardingPass_has_string}
  \uses{def:boarding-pass}
  An accepted boarding pass holds a boarding pass string: its namespace and
  the data element are present, the value is a text string, and the string
  follows Resolution 792, whatever the resolution says.
\end{theorem}
\begin{proof}
  \uses{def:boarding-pass}
  Each of the four ways the string can be wrong is a finding.
\end{proof}

\begin{theorem}[No credential is both a photoid and a boarding pass]
  \label{thm:not-both}
  \lean{Rules.OneId.not_photoid_and_boardingPass}
  \uses{def:photoid, def:boarding-pass, thm:photoid-doctype, thm:bp-doctype}
  No credential is accepted both as a photoid and as a boarding pass.
\end{theorem}
\begin{proof}
  \uses{thm:photoid-doctype, thm:bp-doctype}
  The two document types differ.
\end{proof}

\begin{theorem}[A boarding pass carries only what the guide lists]
  \label{thm:bp-only-listed}
  \lean{Rules.OneId.boardingPass_only_listed}
  \uses{def:boarding-pass}
  An accepted boarding pass carries one namespace with one data element, the
  boarding pass string.  Not proven.

  Counterexample: a boarding pass that carries a second namespace with the
  seat as a data element of its own is accepted under the reading that more
  than the guide lists is allowed.

  Open point: the guide lists one namespace with one data element and says a
  later version may add a namespace.  May a boarding pass credential carry a
  namespace or a data element the guide does not list today?

  Source: \href{https://standards.atlassian.net/wiki/spaces/OneID/pages/2947416082/Technical+Implementation+Guide}{One ID guide} §4.1.2.
\end{theorem}

\begin{definition}[Expiry of a credential (C-44)]
  \label{def:expired}
  \lean{Rules.OneId.expired, Rules.OneId.ExpiryReading}
  \uses{def:mdoc, def:instant}
  A credential is expired when the current time is after the end of the
  validity its issuer signed.  Whether it is still valid at the exact instant
  is a reading: the check is in ISO/IEC 18013-5, which this project has not
  read.

  External: ISO/IEC 18013-5:2021 §9.3.1.

  Source: \href{https://standards.atlassian.net/wiki/spaces/OneID/pages/2947416082/Technical+Implementation+Guide}{One ID guide} §4.1.2, §4.2.1.2.
\end{definition}

\begin{theorem}[Expired after the end of validity]
  \label{thm:expired-after}
  \lean{Rules.OneId.expired_after_validUntil}
  \uses{def:expired}
  A credential is expired once the current time is after the end of its
  validity, under both readings of the exact instant.
\end{theorem}
\begin{proof}
  \uses{def:expired}
  By the definition, in each reading.
\end{proof}

\begin{theorem}[Valid before the end of validity]
  \label{thm:valid-before}
  \lean{Rules.OneId.valid_before_validUntil}
  \uses{def:expired}
  A credential is not expired while the current time is before the end of its
  validity, under both readings of the exact instant.
\end{theorem}
\begin{proof}
  \uses{def:expired}
  By the definition, in each reading.
\end{proof}

\begin{theorem}[An expired credential stays expired]
  \label{thm:expired-stays}
  \lean{Rules.OneId.expired_stays_expired}
  \uses{def:expired}
  Once a credential has expired it is expired at every later time.
\end{theorem}
\begin{proof}
  \uses{def:expired}
  Time only moves further past the end of validity.
\end{proof}

\begin{theorem}[Valid at the end of validity]
  \label{thm:valid-at-validuntil}
  \lean{Rules.OneId.valid_at_validUntil}
  \uses{def:expired}
  A credential is not expired at the exact instant its validity ends.  Not
  proven.

  Counterexample: a boarding pass whose validity ends exactly at departure,
  which the guide allows, is expired at the moment of departure under the
  reading that validity ends at the instant itself.

  Open point: is a credential still valid at the exact instant its validity
  ends?

  Lookup: ISO/IEC 18013-5:2021 §9.3.1.

  Source: \href{https://standards.atlassian.net/wiki/spaces/OneID/pages/2947416082/Technical+Implementation+Guide}{One ID guide} §4.1.2, §4.2.1.2.
\end{theorem}

\begin{definition}[Validity of a boarding pass (C-38)]
  \label{def:bp-validity}
  \lean{Rules.OneId.BoardingPass.validityFindings, Rules.OneId.ValidityReading, Rules.OneId.DepartureReading, Rules.OneId.Journey, Rules.OneId.Journey.departure}
  \uses{def:boarding-pass, def:expired}
  The validity of a boarding pass is accepted only if it ends on or soon
  after the departure of the last leg of the journey.  Which departure is
  meant, the scheduled one or the actual one, and how long ``soon after''
  lasts are a reading: the guide says neither.

  External: ISO/IEC 18013-5:2021, validity information.

  Source: \href{https://standards.atlassian.net/wiki/spaces/OneID/pages/2947416082/Technical+Implementation+Guide}{One ID guide} §4.1.2, implementation consideration 1.
\end{definition}

\begin{theorem}[Validity reaches the departure]
  \label{thm:validity-reaches-departure}
  \lean{Rules.OneId.validity_reaches_departure}
  \uses{def:bp-validity}
  The validity of an accepted boarding pass does not end before the departure
  of the last leg, whichever departure the reading names and however long
  ``soon after'' is.
\end{theorem}
\begin{proof}
  \uses{def:bp-validity}
  An earlier end is a finding.
\end{proof}

\begin{theorem}[Validity ends soon after the departure]
  \label{thm:validity-ends-soon-after}
  \lean{Rules.OneId.validity_ends_soon_after}
  \uses{def:bp-validity}
  The validity of an accepted boarding pass ends no later than ``soon after''
  the departure of the last leg, for the departure and the length the reading
  gives.
\end{theorem}
\begin{proof}
  \uses{def:bp-validity}
  A later end is a finding.
\end{proof}

\begin{theorem}[Valid before the departure]
  \label{thm:valid-before-departure}
  \lean{Rules.OneId.valid_before_departure}
  \uses{def:bp-validity, thm:valid-before, thm:validity-reaches-departure}
  An accepted boarding pass is not expired before the departure of the last
  leg, for the departure the reading names.
\end{theorem}
\begin{proof}
  \uses{thm:valid-before, thm:validity-reaches-departure}
  Its validity reaches the departure, and a credential is valid before the
  end of its validity.
\end{proof}

\begin{theorem}[Expired once the journey is completed]
  \label{thm:expired-once-journey-completed}
  \lean{Rules.OneId.expired_once_journey_completed}
  \uses{def:bp-validity, def:c42}
  An accepted boarding pass is expired once the journey is completed.  Not
  proven.

  Counterexample: a flight leaves at 10:00 and lands at 12:00.  With ``soon
  after'' read as a day, a boarding pass whose validity ends at 10:00 the next
  day is accepted, and at 14:00 it is still valid.

  Open point: the guide says the end of validity makes the boarding pass
  expire once the journey is completed, gives no limit for ``soon after'' and
  no definition of a completed journey, and puts revocation out of scope.  How
  long after the last departure may the validity of a boarding pass run?

  Source: \href{https://standards.atlassian.net/wiki/spaces/OneID/pages/2947416082/Technical+Implementation+Guide}{One ID guide} §4.1.2, implementation consideration 1.
\end{theorem}

\begin{theorem}[Valid until the flight actually departs]
  \label{thm:valid-until-actual-departure}
  \lean{Rules.OneId.valid_until_actual_departure}
  \uses{def:bp-validity}
  An accepted boarding pass is not expired before its last leg actually
  departs.  Not proven.

  Counterexample: a flight scheduled for 10:00 departs at 13:00.  With the
  scheduled departure meant and ``soon after'' read as an hour, a boarding
  pass whose validity ends at 11:00 is accepted, and at 12:30 it is expired:
  the passenger cannot present it to board.

  Open point: an issuer knows only the scheduled departure when it signs.
  Which departure is meant, and what happens to a boarding pass credential
  when its flight is delayed past the end of its validity?

  Source: \href{https://standards.atlassian.net/wiki/spaces/OneID/pages/2947416082/Technical+Implementation+Guide}{One ID guide} §4.1.2, implementation consideration 1.
\end{theorem}

\begin{definition}[Validity of a photoid (C-35)]
  \label{def:photoid-validity}
  \lean{Rules.OneId.Photoid.validityFindings, Rules.OneId.PhotoidValidityReading}
  \uses{def:photoid, def:expired}
  The guide gives no rule for the validity of a photoid.  Two readings are
  kept: the issuer may set any end of validity, or the validity ends no later
  than the document the photoid was derived from expires.

  Source: \href{https://standards.atlassian.net/wiki/spaces/OneID/pages/2947416082/Technical+Implementation+Guide}{One ID guide} §4.1.1, §1.4.
\end{definition}

\begin{theorem}[A photoid expires with its passport]
  \label{thm:photoid-expires-with-document}
  \lean{Rules.OneId.photoid_expires_with_document}
  \uses{def:photoid-validity}
  A photoid whose validity was accepted is expired once the document it was
  derived from has expired.  Not proven.

  Counterexample: a photoid whose validity ends on 1 January 2030, derived
  from a passport that expires on 30 June 2027, is accepted under the reading
  that the issuer may set any end of validity.  Two weeks after the passport
  expired the photoid is still valid.

  Open point: must the validity of a photoid end no later than the expiry of
  the passport it was derived from?

  Source: \href{https://standards.atlassian.net/wiki/spaces/OneID/pages/2947416082/Technical+Implementation+Guide}{One ID guide} §4.1.1, §1.4.
\end{theorem}

\begin{definition}[One traveler for a photoid and a boarding pass (C-41)]
  \label{def:traveler-match}
  \lean{Rules.OneId.travelerMatch, Rules.OneId.MatchReading, Rules.OneId.Issued, Rules.OneId.Traveler, Rules.OneId.Photoid.name}
  \uses{def:photoid, def:boarding-pass}
  A photoid and a boarding pass presented together are accepted as a pair
  only if both are about the same traveler.  The guide states no rule for
  deciding it, so the check is a reading: nothing is checked, which is the
  guide as written, or the names in the two credentials agree.  Other
  readings, such as one that uses the device key, are added as the rules they
  need are modelled.

  External: IATA Resolution 792; ISO/IEC TS 23220-4:2026 Annex C, Table C.1.

  Source: \href{https://standards.atlassian.net/wiki/spaces/OneID/pages/2947416082/Technical+Implementation+Guide}{One ID guide} §4.1.2, §3.2.
\end{definition}

\begin{theorem}[An accepted pair is one traveler's]
  \label{thm:matched-pair-same-traveler}
  \lean{Rules.OneId.matched_pair_same_traveler}
  \uses{def:traveler-match, def:c03, def:c19, def:c42}
  A photoid and a boarding pass that are accepted as a pair are about the
  same traveler.  Not proven.

  Counterexample: with no check between the two, Alice's photoid is accepted
  with Bruno's boarding pass.  With the names compared, the photoid of one
  John Smith is accepted with the boarding pass of another John Smith.

  Open point: the guide requires a boarding pass together with the face image
  of a photoid at a contactless touchpoint and names no data element to
  compare.  No check of the two credentials' content settles it by itself,
  since who a credential is about is fixed at issuance, which the guide puts
  out of scope.  What ties a photoid and a boarding pass to one traveler: a
  rule for issuers, the device key, or something else, and what must a
  relying party check?

  Source: \href{https://standards.atlassian.net/wiki/spaces/OneID/pages/2947416082/Technical+Implementation+Guide}{One ID guide} §4.1.2, §3.2.
\end{theorem}

\begin{theorem}[No check of content alone identifies the traveler]
  \label{thm:no-content-check}
  \lean{Rules.OneId.no_content_check_identifies_the_traveler}
  \uses{def:traveler-match}
  Whatever a relying party checks on the content of a photoid and a boarding
  pass, a pair it accepts can be about two different travelers: the same two
  credentials could have been issued about different people.  So a matching
  rule gives one traveler only together with a rule about issuance.
\end{theorem}
\begin{proof}
  \uses{def:traveler-match}
  The subject of a credential is not part of its content.
\end{proof}

\begin{theorem}[A traveler's own pair is accepted]
  \label{thm:same-traveler-matches}
  \lean{Rules.OneId.same_traveler_matches}
  \uses{def:traveler-match, def:c102}
  A photoid and a boarding pass about the same traveler are accepted as a
  pair.  Not proven.

  Counterexample: a traveler whose passport gives GARCÍA and JOSÉ MARÍA, and
  whose boarding pass carries GARCIA and JOSEMARIA, is refused under an exact
  comparison of names.

  Open point: if names are compared, how are a passport name and a boarding
  pass name brought to the same form?

  Source: \href{https://standards.atlassian.net/wiki/spaces/OneID/pages/2947416082/Technical+Implementation+Guide}{One ID guide} §4.1.2.
\end{theorem}
